\documentclass[10pt,a4paper]{amsart}
\usepackage{amsmath,amssymb,amsthm,graphicx}
% Retained preamble for the Stacks Project chapter snapshot in this directory.
% The chapter's own preamble.tex is a copy of the shared Stacks Project preamble; the block
% below is sliced MECHANICALLY, line for line, from it. Nothing was added, removed or
% redefined; the wrapper class, the xr-hyper/hyperref wiring and the externaldocument list
% are omitted because the assembler supplies the document class.
%
% lines taken from preamble.tex: 3, 5, 9, 11, 12, 156, 157, 158, ... (51 lines total)
%
% package coverage restored from the chapter's own preamble (2026-10-13): xy, xy-2cell

\usepackage[all]{xy}
\xyoption{2cell}
\UseAllTwocells

\newenvironment{reference}{\comment}{\endcomment}
\newenvironment{slogan}{\comment}{\endcomment}
\newenvironment{history}{\comment}{\endcomment}
\theoremstyle{plain}
\newtheorem{theorem}[subsection]{Theorem}
\newtheorem{proposition}[subsection]{Proposition}
\newtheorem{lemma}[subsection]{Lemma}
\theoremstyle{definition}
\newtheorem{definition}[subsection]{Definition}
\newtheorem{example}[subsection]{Example}
\newtheorem{exercise}[subsection]{Exercise}
\newtheorem{situation}[subsection]{Situation}
\theoremstyle{remark}
\newtheorem{remark}[subsection]{Remark}
\newtheorem{remarks}[subsection]{Remarks}
\numberwithin{equation}{subsection}
\def\lim{\mathop{\mathrm{lim}}\nolimits}
\def\colim{\mathop{\mathrm{colim}}\nolimits}
\def\Spec{\mathop{\mathrm{Spec}}}
\def\Hom{\mathop{\mathrm{Hom}}\nolimits}
\def\Ext{\mathop{\mathrm{Ext}}\nolimits}
\def\SheafHom{\mathop{\mathcal{H}\!\mathit{om}}\nolimits}
\def\SheafExt{\mathop{\mathcal{E}\!\mathit{xt}}\nolimits}
\def\Sch{\mathit{Sch}}
\def\Mor{\mathop{\mathrm{Mor}}\nolimits}
\def\Ob{\mathop{\mathrm{Ob}}\nolimits}
\def\Sh{\mathop{\mathit{Sh}}\nolimits}
\def\NL{\mathop{N\!L}\nolimits}
\def\CH{\mathop{\mathrm{CH}}\nolimits}
\def\proetale{{pro\text{-}\acute{e}tale}}
\def\etale{{\acute{e}tale}}
\def\QCoh{\mathit{QCoh}}
\def\Ker{\mathop{\mathrm{Ker}}}
\def\Im{\mathop{\mathrm{Im}}}
\def\Coker{\mathop{\mathrm{Coker}}}
\def\Coim{\mathop{\mathrm{Coim}}}
\DeclareMathSymbol{\boxtimes}{\mathbin}{AMSa}{"02}
\def\QCohstack{\mathcal{QC}\!\mathit{oh}}
\def\Cohstack{\mathcal{C}\!\mathit{oh}}
\def\Spacesstack{\mathcal{S}\!\mathit{paces}}
\def\Quotfunctor{\mathrm{Quot}}
\def\Hilbfunctor{\mathrm{Hilb}}
\def\Curvesstack{\mathcal{C}\!\mathit{urves}}
\def\Polarizedstack{\mathcal{P}\!\mathit{olarized}}
\def\Complexesstack{\mathcal{C}\!\mathit{omplexes}}
\def\Pic{\mathop{\mathrm{Pic}}\nolimits}
\def\Picardstack{\mathcal{P}\!\mathit{ic}}
\def\Picardfunctor{\mathrm{Pic}}
\def\Deformationcategory{\mathcal{D}\!\mathit{ef}}
\usepackage{cleveref}
\crefname{originallemma}{original lemma}{original lemmas}
\Crefname{originallemma}{Original lemma}{Original lemmas}
\makeatletter
\newlabel{equation-to-solve-ringed-spaces}{{08U7}{source}{}{sourcepointer-1}{}}
\newlabel{equation-to-solve-ringed-spaces@cref}{{[originallemma][0][]08U7}{[1][0][]source}{}{}{}}
\makeatother
\crefname{originallemma}{original lemma}{original lemmas}
\Crefname{originallemma}{Original lemma}{Original lemmas}
\makeatletter
\newlabel{lemma-huge-diagram}{{08S5}{source}{}{sourcepointer-2}{}}
\newlabel{lemma-huge-diagram@cref}{{[originallemma][0][]08S5}{[1][0][]source}{}{}{}}
\makeatother
\crefname{originallemma}{original lemma}{original lemmas}
\Crefname{originallemma}{Original lemma}{Original lemmas}
\makeatletter
\newlabel{modules-definition-cotangent-complex-morphism-ringed-topoi}{{08TN}{source}{}{sourcepointer-3}{}}
\newlabel{modules-definition-cotangent-complex-morphism-ringed-topoi@cref}{{[originallemma][0][]08TN}{[1][0][]source}{}{}{}}
\makeatother
\crefname{originallemma}{original lemma}{original lemmas}
\Crefname{originallemma}{Original lemma}{Original lemmas}
\makeatletter
\newlabel{modules-lemma-NL-up-to-qis}{{08TL}{source}{}{sourcepointer-4}{}}
\newlabel{modules-lemma-NL-up-to-qis@cref}{{[originallemma][0][]08TL}{[1][0][]source}{}{}{}}
\makeatother
\crefname{originallemma}{original lemma}{original lemmas}
\Crefname{originallemma}{Original lemma}{Original lemmas}
\makeatletter
\newlabel{modules-lemma-pullback-NL}{{08TM}{source}{}{sourcepointer-5}{}}
\newlabel{modules-lemma-pullback-NL@cref}{{[originallemma][0][]08TM}{[1][0][]source}{}{}{}}
\makeatother
\title{Stacks, Deformation Theory: lemma-huge-diagram-ringed-spaces}
\author{The Stacks Project Authors}
\date{stacks-project:defos.tex:local-snapshot-20261012; GNU FDL 1.2 or later}
\begin{document}
\maketitle
\noindent\textit{Editorial note.} Statement span is lines 1955-2013 of sources/stacks-defos/pinned-original-source.tex; proof span is lines 2015-2155, the complete author proof (139 body lines).\par\medskip
\noindent\textit{Editorial note.} Navigation stubs: 5 cited-but-not-included labels are supplied as non-load-bearing pointers with their pinned tags where the snapshot tag index has them.\par\medskip
\noindent\textit{Editorial note.} extra\_packages is empty on purpose: the retained preamble already loads xy with options and a second declaration clashes in the assembler wrapper.\par\medskip
\noindent\textit{Editorial note.} Text is retained verbatim, including the author's own typos and placeholders.\par\medskip
\noindent\textit{Editorial note.} Queue index 1 of queue\_defos.json (pb=5962, stmt 1955-2013, proof 2015-2155).\par\medskip
\noindent\textit{Allowed standard prerequisites (editorial).}\begin{itemize}
\item Explicit prerequisite boundary. The cited inputs are prerequisites; the argument in the retained span is not replaced by any cited result.
\end{itemize}\medskip
\ifdefined\hypertarget\hypertarget{sourcepointer-1}{}\fi
\noindent\textit{Source reference.} Original lemma 08U7: sources/stacks-defos/pinned-original-source.tex, source lines 1935--1943. Cited by the retained statement or proof but not part of the collected spans; supplied only as a navigation pointer. The argument does not depend on its content.\par\medskip
\ifdefined\hypertarget\hypertarget{sourcepointer-2}{}\fi
\noindent\textit{Source reference.} Original lemma 08S5: sources/stacks-defos/pinned-original-source.tex, source lines 275--296. Cited by the retained statement or proof but not part of the collected spans; supplied only as a navigation pointer. The argument does not depend on its content.\par\medskip
\ifdefined\hypertarget\hypertarget{sourcepointer-3}{}\fi
\noindent\textit{Source reference.} Original lemma 08TN: sources/stacks-defos/pinned-original-source.tex, source lines 2155--2155. Cited by the retained statement or proof but not part of the collected spans; supplied only as a navigation pointer. The argument does not depend on its content.\par\medskip
\ifdefined\hypertarget\hypertarget{sourcepointer-4}{}\fi
\noindent\textit{Source reference.} Original lemma 08TL: sources/stacks-defos/pinned-original-source.tex, source lines 2155--2155. Cited by the retained statement or proof but not part of the collected spans; supplied only as a navigation pointer. The argument does not depend on its content.\par\medskip
\ifdefined\hypertarget\hypertarget{sourcepointer-5}{}\fi
\noindent\textit{Source reference.} Original lemma 08TM: sources/stacks-defos/pinned-original-source.tex, source lines 2155--2155. Cited by the retained statement or proof but not part of the collected spans; supplied only as a navigation pointer. The argument does not depend on its content.\par\medskip
\begin{lemma}
\label{lemma-huge-diagram-ringed-spaces}
Assume given a commutative diagram of morphisms of ringed spaces
\begin{equation}
\label{equation-huge-1}
\vcenter{
\xymatrix{
& (X_2, \mathcal{O}_{X_2}) \ar[r]_{i_2} \ar[d]_{f_2} \ar[ddl]_g &
(X'_2, \mathcal{O}_{X'_2}) \ar[d]^{f'_2} \\
& (S_2, \mathcal{O}_{S_2}) \ar[r]^{t_2} \ar[ddl]|\hole &
(S'_2, \mathcal{O}_{S'_2}) \ar[ddl] \\
(X_1, \mathcal{O}_{X_1}) \ar[r]_{i_1} \ar[d]_{f_1} &
(X'_1, \mathcal{O}_{X'_1}) \ar[d]^{f'_1} \\
(S_1, \mathcal{O}_{S_1}) \ar[r]^{t_1} &
(S'_1, \mathcal{O}_{S'_1})
}
}
\end{equation}
whose horizontal arrows are first order thickenings. Set
$\mathcal{G}_j = \Ker(i_j^\sharp)$ and assume given
a $g$-map $\nu : \mathcal{G}_1 \to \mathcal{G}_2$ of modules
giving rise to the commutative diagram
\begin{equation}
\label{equation-huge-2}
\vcenter{
\xymatrix{
& 0 \ar[r] & \mathcal{G}_2 \ar[r] &
\mathcal{O}_{X'_2} \ar[r] &
\mathcal{O}_{X_2} \ar[r] & 0 \\
& 0 \ar[r]|\hole &
\mathcal{J}_2 \ar[u]_{c_2} \ar[r] &
\mathcal{O}_{S'_2} \ar[u] \ar[r]|\hole &
\mathcal{O}_{S_2} \ar[u] \ar[r] & 0 \\
0 \ar[r] & \mathcal{G}_1 \ar[ruu] \ar[r] &
\mathcal{O}_{X'_1} \ar[r] &
\mathcal{O}_{X_1} \ar[ruu] \ar[r] & 0 \\
0 \ar[r] & \mathcal{J}_1 \ar[ruu]|\hole \ar[u]^{c_1} \ar[r] &
\mathcal{O}_{S'_1} \ar[ruu]|\hole \ar[u] \ar[r] &
\mathcal{O}_{S_1} \ar[ruu]|\hole \ar[u] \ar[r] & 0
}
}
\end{equation}
with front and back solutions to (\ref{equation-to-solve-ringed-spaces}).
\begin{enumerate}
\item There exist a canonical element in
$\Ext^1_{\mathcal{O}_{X_2}}(Lg^*\NL_{X_1/S_1}, \mathcal{G}_2)$
whose vanishing is a necessary and sufficient condition for the existence
of a morphism of ringed spaces $X'_2 \to X'_1$ fitting into
(\ref{equation-huge-1}) compatibly with $\nu$.
\item If there exists a morphism $X'_2 \to X'_1$ fitting into
(\ref{equation-huge-1}) compatibly with $\nu$ the set of all such morphisms
is a principal homogeneous space under
$$
\Hom_{\mathcal{O}_{X_1}}(\Omega_{X_1/S_1}, g_*\mathcal{G}_2) =
\Hom_{\mathcal{O}_{X_2}}(g^*\Omega_{X_1/S_1}, \mathcal{G}_2) =
\Ext^0_{\mathcal{O}_{X_2}}(Lg^*\NL_{X_1/S_1}, \mathcal{G}_2).
$$
\end{enumerate}
\end{lemma}

\begin{proof}
The naive cotangent complex $\NL_{X_1/S_1}$ is defined in Modules, Definition
\ref{modules-definition-cotangent-complex-morphism-ringed-topoi}.
The equalities in the last statement of the lemma follow from
the fact that $g^*$ is adjoint to $g_*$, the fact that
$H^0(\NL_{X_1/S_1}) = \Omega_{X_1/S_1}$ (by construction of the
naive cotangent complex) and the fact that $Lg^*$ is the left
derived functor of $g^*$. Thus we will work with the groups
$\Ext^k_{\mathcal{O}_{X_2}}(Lg^*\NL_{X_1/S_1}, \mathcal{G}_2)$,
$k = 0, 1$ in the rest of the proof. We first argue that we can reduce
to the case where the underlying topological spaces of all ringed
spaces in the lemma is the same.

\medskip\noindent
To do this, observe that $g^{-1}\NL_{X_1/S_1}$ is equal to the naive
cotangent complex of the homomorphism of sheaves of rings
$g^{-1}f_1^{-1}\mathcal{O}_{S_1} \to g^{-1}\mathcal{O}_{X_1}$, see
Modules, Lemma \ref{modules-lemma-pullback-NL}.
Moreover, the degree $0$ term of $\NL_{X_1/S_1}$ is a flat
$\mathcal{O}_{X_1}$-module, hence the canonical map
$$
Lg^*\NL_{X_1/S_1}
\longrightarrow
g^{-1}\NL_{X_1/S_1} \otimes_{g^{-1}\mathcal{O}_{X_1}} \mathcal{O}_{X_2}
$$
induces an isomorphism on cohomology sheaves in degrees $0$ and $-1$.
Thus we may replace the Ext groups of the lemma with
$$
\Ext^k_{g^{-1}\mathcal{O}_{X_1}}(g^{-1}\NL_{X_1/S_1}, \mathcal{G}_2) =
\Ext^k_{g^{-1}\mathcal{O}_{X_1}}(
\NL_{g^{-1}\mathcal{O}_{X_1}/g^{-1}f_1^{-1}\mathcal{O}_{S_1}}, \mathcal{G}_2)
$$
The set of morphism of ringed spaces $X'_2 \to X'_1$ fitting into
(\ref{equation-huge-1}) compatibly with $\nu$
is in one-to-one bijection with
the set of homomorphisms of $g^{-1}f_1^{-1}\mathcal{O}_{S'_1}$-algebras
$g^{-1}\mathcal{O}_{X'_1} \to \mathcal{O}_{X'_2}$ which are compatible with
$f^\sharp$ and $\nu$. In this way we see that we may assume we have a
diagram (\ref{equation-huge-2}) of sheaves on $X$ and we are looking to
find a homomorphism of sheaves of rings
$\mathcal{O}_{X'_1} \to \mathcal{O}_{X'_2}$ fitting into it.

\medskip\noindent
In the rest of the proof of the lemma we assume
all underlying topological spaces are the
same, i.e., we have a diagram (\ref{equation-huge-2}) of sheaves on
a space $X$ and we are looking for homomorphisms of sheaves of rings
$\mathcal{O}_{X'_1} \to \mathcal{O}_{X'_2}$ fitting into it.
As ext groups we will use
$\Ext^k_{\mathcal{O}_{X_1}}(
\NL_{\mathcal{O}_{X_1}/\mathcal{O}_{S_1}}, \mathcal{G}_2)$, $k = 0, 1$.

\medskip\noindent
Step 1. Construction of the obstruction class. Consider the sheaf
of sets
$$
\mathcal{E} = \mathcal{O}_{X'_1} \times_{\mathcal{O}_{X_2}} \mathcal{O}_{X'_2}
$$
This comes with a surjective map $\alpha : \mathcal{E} \to \mathcal{O}_{X_1}$
and hence we can use $\NL(\alpha)$ instead of
$\NL_{\mathcal{O}_{X_1}/\mathcal{O}_{S_1}}$, see
Modules, Lemma \ref{modules-lemma-NL-up-to-qis}.
Set
$$
\mathcal{I}' =
\Ker(\mathcal{O}_{S'_1}[\mathcal{E}] \to \mathcal{O}_{X_1})
\quad\text{and}\quad
\mathcal{I} =
\Ker(\mathcal{O}_{S_1}[\mathcal{E}] \to \mathcal{O}_{X_1})
$$
There is a surjection $\mathcal{I}' \to \mathcal{I}$ whose kernel
is $\mathcal{J}_1\mathcal{O}_{S'_1}[\mathcal{E}]$.
We obtain two homomorphisms of $\mathcal{O}_{S'_1}$-algebras
$$
a : \mathcal{O}_{S'_1}[\mathcal{E}] \to \mathcal{O}_{X'_1}
\quad\text{and}\quad
b : \mathcal{O}_{S'_1}[\mathcal{E}] \to \mathcal{O}_{X'_2}
$$
which induce maps $a|_{\mathcal{I}'} : \mathcal{I}' \to \mathcal{G}_1$ and
$b|_{\mathcal{I}'} : \mathcal{I}' \to \mathcal{G}_2$. Both $a$ and $b$
annihilate $(\mathcal{I}')^2$. Moreover $a$ and $b$ agree on
$\mathcal{J}_1\mathcal{O}_{S'_1}[\mathcal{E}]$ as maps into $\mathcal{G}_2$
because the left hand square of (\ref{equation-huge-2}) is commutative.
Thus the difference
$b|_{\mathcal{I}'} - \nu \circ a|_{\mathcal{I}'}$
induces a well defined $\mathcal{O}_{X_1}$-linear map
$$
\xi : \mathcal{I}/\mathcal{I}^2 \longrightarrow \mathcal{G}_2
$$
which sends the class of a local section $f$ of $\mathcal{I}$ to
$\nu(a(f')) - b(f')$ where $f'$ is a lift of $f$ to a local
section of $\mathcal{I}'$. We let
$[\xi] \in \Ext^1_{\mathcal{O}_{X_1}}(\NL(\alpha), \mathcal{G}_2)$
be the image (see below).

\medskip\noindent
Step 2. Vanishing of $[\xi]$ is necessary. Let us write
$\Omega = \Omega_{\mathcal{O}_{S_1}[\mathcal{E}]/\mathcal{O}_{S_1}}
\otimes_{\mathcal{O}_{S_1}[\mathcal{E}]} \mathcal{O}_{X_1}$.
Observe that $\NL(\alpha) = (\mathcal{I}/\mathcal{I}^2 \to \Omega)$
fits into a distinguished triangle
$$
\Omega[0] \to
\NL(\alpha) \to
\mathcal{I}/\mathcal{I}^2[1] \to
\Omega[1]
$$
Thus we see that $[\xi]$ is zero if and only if $\xi$
is a composition $\mathcal{I}/\mathcal{I}^2 \to \Omega \to \mathcal{G}_2$
for some map $\Omega \to \mathcal{G}_2$. Suppose there exists a
homomorphisms of sheaves of rings
$\varphi : \mathcal{O}_{X'_1} \to \mathcal{O}_{X'_2}$ fitting into
(\ref{equation-huge-2}). In this case consider the map
$\mathcal{O}_{S'_1}[\mathcal{E}] \to \mathcal{G}_2$,
$f' \mapsto b(f') - \varphi(a(f'))$. A calculation
shows this annihilates $\mathcal{J}_1\mathcal{O}_{S'_1}[\mathcal{E}]$
and induces a derivation $\mathcal{O}_{S_1}[\mathcal{E}] \to \mathcal{G}_2$.
The resulting linear map $\Omega \to \mathcal{G}_2$ witnesses the
fact that $[\xi] = 0$ in this case.

\medskip\noindent
Step 3. Vanishing of $[\xi]$ is sufficient. Let
$\theta : \Omega \to \mathcal{G}_2$ be a $\mathcal{O}_{X_1}$-linear map
such that $\xi$ is equal to
$\theta \circ (\mathcal{I}/\mathcal{I}^2 \to \Omega)$.
Then a calculation shows that
$$
b + \theta \circ d : \mathcal{O}_{S'_1}[\mathcal{E}] \to \mathcal{O}_{X'_2}
$$
restricted to $\mathcal{I}'$ agrees with
$\nu \circ a : \mathcal{I}' \to \mathcal{G}_2$. Since $\mathcal{O}_{X'_1}$
is the pushout of $\mathcal{I}' \to \mathcal{O}_{S'_1}[\mathcal{E}]$
and $\mathcal{I}' \to \mathcal{G}_1$ the maps $b + \theta \circ d$
and $a$ define a map
$\mathcal{O}_{X'_1} \to \mathcal{O}_{X'_2}$ fitting into
(\ref{equation-huge-2}).

\medskip\noindent
Proof of (2) in the special case above. Omitted. Hint:
This is exactly the same as the proof of (2) of Lemma \ref{lemma-huge-diagram}.
\end{proof}

\end{document}
